\section{Overview}\label{2.1}

This chapter surveys the four bodies of prior work that AURA-Bio's design
draws on directly: retrieval-augmented generation, hallucination in
language models, dense and sparse retrieval, and the case against using a
second language model as the sole judge of a first one's output. It closes
with a short account of intelligent tutoring systems, since AURA-Bio's
adaptive-teaching component is not a new idea in education, only a new
setting for it.

\section{Retrieval-Augmented Generation}\label{2.2}

The idea of pairing a language model with an external, searchable knowledge
source rather than relying purely on what the model memorised during
training was formalised by Lewis et al.\ \cite{lewis2020rag}, whose original
Retrieval-Augmented Generation (RAG) architecture retrieves a small number
of relevant passages and conditions the model's output on them. The appeal
is straightforward: a model can be kept current and be made to answer
about specific documents without retraining it, simply by changing what it
is allowed to read at answer time. Gao et al.\ \cite{gao2023ragsurvey}
survey how the field has grown since, through better chunking, better
retrieval, and multi-stage pipelines. But the surveyed systems share a
structural limitation directly relevant to this project: most stop once an
answer has been generated from the retrieved context. They do not
independently check whether the generated answer actually agrees with what
it cites, nor whether it omitted something the underlying model already
knew. Both gaps are addressed in AURA-Bio's verification layer, described in
Chapter~3, and both turned out in practice to matter (Chapter~4).

Shuster et al.\ \cite{shuster2021retrieval} showed specifically that
retrieval augmentation reduces, but does not eliminate hallucination in
conversational models, which motivates treating retrieval as one layer of
defence rather than a complete solution on its own.

\section{Hallucination in Language Models}\label{2.3}

Ji et al.\ \cite{ji2023hallucination} provide the standard survey of
hallucination in natural-language generation that describes model output that is fluent
and plausible but not grounded in any real source or fact. Their taxonomy
distinguishes \emph{intrinsic} hallucination (contradicting the given
source) from \emph{extrinsic} hallucination (asserting something the source
neither confirms nor denies) which is a distinction that maps directly onto
AURA-Bio's claim-support labels of \texttt{contradictory} versus
\texttt{unsupported} versus \texttt{unverifiable} (Chapter~3). A recurring
theme in this literature, and one this project's own measurements
independently confirmed (Chapter~4), is that hallucination rates reported
by an automated benchmark are only as trustworthy as the scorer that
produces them: a scorer that looks for a fixed set of literal phrases will
systematically miscount a model that expresses the same idea in different
words.

\section{Dense and Sparse Retrieval}\label{2.4}

Sparse retrieval based on term matching has a long history, and the BM25
ranking function described by Robertson and Zaragoza
\cite{robertson2009bm25} remains, decades later, a strong and
computationally cheap baseline, particularly for exact terminology:
gene names, enzyme names, and abbreviations, where a dense embedding can
blur distinctions that a keyword match preserves exactly. Dense retrieval,
by contrast, represents text as vectors positioned so that semantically
similar passages sit close together regardless of exact wording. Reimers
and Gurevych's Sentence-BERT \cite{reimers2019sbert} made it practical to
produce high-quality sentence embeddings efficiently, and the BGE family of
embedding models \cite{xiao2023bge} used in this project builds on that
line of work with training specifically aimed at retrieval quality.
Domain-specific embeddings matter too: Liu et al.'s SapBERT
\cite{liu2021sapbert} demonstrated that an encoder pretrained on biomedical
entity synonyms captures distinctions a general-purpose encoder misses,
which is why a biomedical encoder is one of the components evaluated for
this project's learned question router (Chapter~3). Fast approximate
nearest-neighbour search over millions of such vectors is made practical by
FAISS \cite{johnson2019faiss}, which underpins the dense index built for
this project.

Neither channel alone is sufficient: sparse retrieval misses paraphrase and
synonym, dense retrieval can miss exact rare terms. Combining the two,
reciprocal rank fusion, in this project's case, and then re-scoring the
combined candidates with a cross-encoder reranker is now standard practice
in retrieval-heavy systems, and is the retrieval architecture adopted here.

\section{Why Not Use a Language Model as the Judge}\label{2.5}

A tempting shortcut for evaluating whether an answer is correct is to ask
another, larger language model to grade it. Zheng et al.\
\cite{zheng2023judging} studied this practice carefully and documented real
biases in it. A preference for longer answers regardless of quality, a
tendency to prefer answers that resemble the judge model's own style, and
inconsistency between repeated judgements of the same answer. For a
tutoring system whose entire purpose is to be more trustworthy than an
unsupervised model, grading its own homework with a second model would
undermine the premise. AURA-Bio's benchmark (Chapter~3) is therefore scored
by machine-checkable answer keys including required concepts, numerical answers
recomputed independently, and canonical process orderings, so that a
result can be verified by inspection rather than trusted on the strength of
another model's opinion. The internal verification layer is deterministic
for the same reason that arithmetic is recomputed with SymPy
\cite{meurer2017sympy} rather than read off the model's own working, and
sequence claims are checked against the standard genetic code with
Biopython \cite{cock2009biopython} rather than judged by a second model's
impression of whether they sound right.

\section{Parameter-Efficient Fine-Tuning}\label{2.6}

Low-Rank Adaptation, or LoRA \cite{hu2021lora}, freezes a pretrained
model's weights and learns a small pair of low-rank update matrices
instead, cutting the number of trainable parameters and hence, the memory
needed to train them, by orders of magnitude. QLoRA
\cite{dettmers2023qlora} extended this further by quantising the frozen
base model to 4-bit precision during training, making it feasible to
fine-tune a several-billion-parameter model on a single consumer GPU with a
few gigabytes of memory. This is the technique planned for AURA-Bio's
fine-tuning phase (Chapter~5): full fine-tuning of even a small model is
outside the reach of the project's hardware, but a QLoRA adapter is not.

\section{The Base Model}\label{2.7}

The local model driving AURA-Bio's generation is a member of the Qwen
family, \textbf{Qwen3.5:4b}, chosen for its combination of small size,
strong instruction-following, and a permissive licence suitable for local,
offline deployment, the same properties the encoder family
\cite{devlin2019bert,reimers2019sbert} was chosen for on the retrieval
side.

\section{Intelligent Tutoring Systems}\label{2.8}

Adaptive teaching did not begin with language models. VanLehn's review of
intelligent tutoring systems \cite{vanlehn2011tutoring} found that
well-designed step-based tutoring systems can approach the effectiveness of
human tutoring, and identified adapting to the learner's current level and
giving targeted, specific feedback rather than generic praise or
correction, as the features that separate an effective tutor from an
ineffective one. AURA-Bio's teaching-style router and learner-level
adaptation (Chapter~3) are a direct, if modernised, continuation of that
finding: the mechanism for adapting has changed from scripted branching
logic to a routed language-model prompt, but the underlying pedagogical
principle, '\emph{what} is explained and \emph{how} it is explained
should depend on who is asking' is unchanged.

\section{Summary}\label{2.9}

\begin{table}[h]
\centering
\caption{Summary of literature reviewed and its relation to AURA-Bio}
\label{tab:2.1}
\footnotesize
\setlength{\tabcolsep}{4pt}
\renewcommand{\arraystretch}{1.2}
\begin{tabularx}{\textwidth}{@{}p{2.6cm} X X@{}}
\toprule
\textbf{Area} & \textbf{Key reference(s)} & \textbf{Relation to this project} \\
\midrule
Retrieval-augmented generation & \cite{lewis2020rag,gao2023ragsurvey} &
Baseline architecture. AURA-Bio adds independent post-hoc verification that
most surveyed RAG systems lack. \\
\addlinespace[3pt]
Hallucination & \cite{ji2023hallucination,shuster2021retrieval} &
Motivates treating retrieval as one layer, not a complete fix. The
intrinsic/extrinsic distinction maps onto the project's claim labels. \\
\addlinespace[3pt]
Sparse retrieval & \cite{robertson2009bm25} &
BM25 channel for exact terminology, fused with dense retrieval. \\
\addlinespace[3pt]
Dense retrieval \& embeddings & \cite{reimers2019sbert,xiao2023bge,liu2021sapbert,johnson2019faiss} &
Embedding model and index technology for semantic retrieval and for the
learned question router. \\
\addlinespace[3pt]
LLM-as-judge risk & \cite{zheng2023judging} &
Justifies the project's machine-checkable, deterministic benchmark rather
than a second-model judge. \\
\addlinespace[3pt]
Deterministic verification tools & \cite{meurer2017sympy,cock2009biopython} &
SymPy and Biopython used directly inside the verification layer. \\
\addlinespace[3pt]
Parameter-efficient fine-tuning & \cite{hu2021lora,dettmers2023qlora} &
Planned method for the in-progress fine-tuning phase (Chapter~5). \\
\addlinespace[3pt]
Base model family & \cite{qwen2023technical} &
The local generator AURA-Bio wraps. \\
\addlinespace[3pt]
Intelligent tutoring systems & \cite{vanlehn2011tutoring} &
Pedagogical grounding for the teaching-style and learner-level router. \\
\bottomrule
\end{tabularx}
\end{table}
